% This document was generated by an LLM.

\documentclass{worksheet}

\DeclareMathOperator{\arccsc}{arccsc}
\DeclareMathOperator{\arcsec}{arcsec}
\DeclareMathOperator{\arccot}{arccot}
\DeclareMathOperator{\sgn}{sgn}

\title{Trigonometric Substitutions and Inverse Trigonometric Expressions}
\subtitle{Recovering algebraic expressions in $x$}
\worksheetrunninghead{Trigonometric substitution and inverse-trig practice}

\begin{document}
\maketitle

\ifshowsolutions
This version contains worked solutions. Square roots denote the nonnegative
principal square root, and every sign choice follows from the stated range of
$\theta$.
\else
For every problem, express the requested quantity entirely in terms of $x$.
Use the stated substitution or introduce an angle such as
$\theta=\arcsin u$, simplify radicals, choose signs from the relevant range,
and state the domain when it is not already stated.
\fi

\section*{Reference sheet}

\[
\sin^2\theta+\cos^2\theta=1,
\qquad 1+\tan^2\theta=\sec^2\theta,
\qquad 1+\cot^2\theta=\csc^2\theta.
\]
\[
\sin(2\theta)=2\sin\theta\cos\theta,
\quad
\cos(2\theta)=1-2\sin^2\theta=2\cos^2\theta-1,
\quad
\tan(2\theta)=\frac{2\tan\theta}{1-\tan^2\theta}.
\]

Principal ranges used below are
\[
\arcsin u\in[-\pi/2,\pi/2],\qquad
\arccos u\in[0,\pi],\qquad
\arctan u\in(-\pi/2,\pi/2).
\]
For $\arcsec$ problems with positive argument restrictions, use
$\theta\in[0,\pi/2)$. For unrestricted $\arcsec$, use
$\theta\in[0,\pi]$ with $\sin\theta\geq 0$. For $\arccsc$, use
$\theta\in[-\pi/2,\pi/2]$ with $\cos\theta\geq 0$, and for $\arccot$, use
$\theta\in(0,\pi)$.

\section{Direct Substitution Recovery}

Use the stated substitution and angle range. These are the conversions that
often occur at the final stage of a trigonometric-substitution integral.

\begin{enumerate}
\exercise{$x=3\sin\theta$, where $-\pi/2\leq\theta\leq\pi/2$.
Express $2\cos\theta$ in terms of $x$.}{34mm}{%
From $\sin\theta=x/3$,
\[
\cos\theta=\sqrt{1-\sin^2\theta}
=\sqrt{1-\frac{x^2}{9}}=\frac{\sqrt{9-x^2}}{3}.
\]
Cosine is nonnegative on the stated range, so
\[
\boxed{2\cos\theta=\frac{2\sqrt{9-x^2}}{3}},\qquad |x|\leq3.
\]}

\exercise{$x=5\tan\theta$, where $-\pi/2<\theta<\pi/2$.
Express $\tfrac12\sin\theta$ in terms of $x$.}{36mm}{%
Here $\tan\theta=x/5$ and $\sec\theta>0$. Thus
\[
\sec\theta=\sqrt{1+\tan^2\theta}=\frac{\sqrt{x^2+25}}{5}.
\]
Since $\sin\theta=\tan\theta/\sec\theta$,
\[
\boxed{\frac12\sin\theta=\frac{x}{2\sqrt{x^2+25}}}.
\]}

\exercise{$x=2\sec\theta$, where $0\leq\theta<\pi/2$.
Express $\tfrac14\sin\theta$ in terms of $x$.}{36mm}{%
We have $\cos\theta=2/x$. Since sine is nonnegative,
\[
\sin\theta=\sqrt{1-\frac{4}{x^2}}
=\frac{\sqrt{x^2-4}}{x}.
\]
Therefore
\[
\boxed{\frac14\sin\theta=\frac{\sqrt{x^2-4}}{4x}},\qquad x\geq2.
\]}

\exercise{$x=4\sin\theta$, where $-\pi/2<\theta<\pi/2$.
Express $\tan\theta$ in terms of $x$.}{34mm}{%
We have $\sin\theta=x/4$ and
$\cos\theta=\sqrt{16-x^2}/4>0$. Therefore
\[
\boxed{\tan\theta=\frac{x}{\sqrt{16-x^2}}},\qquad |x|<4.
\]}

\exercise{$x=3\tan\theta$, where $-\pi/2<\theta<\pi/2$.
Express $\sec\theta$ in terms of $x$.}{32mm}{%
Because $\sec\theta>0$,
\[
\sec\theta=\sqrt{1+\tan^2\theta}
=\sqrt{1+\frac{x^2}{9}}.
\]
Hence
\[
\boxed{\sec\theta=\frac{\sqrt{x^2+9}}{3}}.
\]}

\exercise{$x=6\sec\theta$, where $0\leq\theta<\pi/2$.
Express $\tan\theta$ in terms of $x$.}{32mm}{%
Since $\sec\theta=x/6$ and $\tan\theta\geq0$,
\[
\boxed{\tan\theta=\sqrt{\sec^2\theta-1}
=\frac{\sqrt{x^2-36}}{6}},\qquad x\geq6.
\]}

\exercise{$x=2\sin\theta$, where $0<\theta\leq\pi/2$.
Express $3\csc\theta$ in terms of $x$.}{28mm}{%
Since $\sin\theta=x/2$,
\[
\boxed{3\csc\theta=\frac{3}{\sin\theta}=\frac{6}{x}},
\qquad 0<x\leq2.
\]}

\exercise{$x=7\tan\theta$, where $-\pi/2<\theta<\pi/2$.
Express $4\cos\theta$ in terms of $x$.}{32mm}{%
We have $\sec\theta=\sqrt{x^2+49}/7$. Therefore
\[
\boxed{4\cos\theta=\frac{4}{\sec\theta}
=\frac{28}{\sqrt{x^2+49}}}.
\]}

\exercise{$x=5\sec\theta$, where $0<\theta<\pi/2$.
Express $2\cot\theta$ in terms of $x$.}{32mm}{%
The identity $\tan^2\theta=\sec^2\theta-1$ gives
$\tan\theta=\sqrt{x^2-25}/5$. Hence
\[
\boxed{2\cot\theta=\frac{10}{\sqrt{x^2-25}}},\qquad x>5.
\]}
\end{enumerate}

\section{Reciprocal Compositions}

Start by setting $\theta$ equal to the inverse trigonometric expression.
These problems require little more than taking a reciprocal, but domains still
matter.

\begin{enumerate}[resume]
\exercise{$5\csc\!\left(\arcsin\dfrac{x}{5}\right)$}{26mm}{%
$\sin\theta = x/5$, so $\csc\theta = 5/x$ and the expression is
$\dfrac{25}{x}$ for $0<|x|\leq 5$.}

\exercise{$\sec\!\left(\arccos 3x\right)$}{24mm}{%
$\cos\theta = 3x$, so the value is $\dfrac{1}{3x}$ for
$0<|x|\leq 1/3$.}

\exercise{$2\cot\!\left(\arctan 2x\right)$}{24mm}{%
$\tan\theta = 2x$, so $\cot\theta = 1/(2x)$ and the value is
$\dfrac{1}{x}$ for $x\neq 0$.}

\exercise{$\cos\!\left(\arcsec\dfrac{x}{4}\right)$}{24mm}{%
$\sec\theta = x/4$, so the value is $\dfrac{4}{x}$ for $|x|\geq 4$.}

\exercise{$\tan\!\left(\arccot 5x\right)$}{30mm}{%
$\cot\theta = 5x$, giving $\dfrac{1}{5x}$ for $x\neq 0$. At $x=0$ we have
$\theta=\pi/2$ and the tangent is undefined, so the point must be excluded
even though $\arccot$ is defined there.}
\end{enumerate}

\section{One Pythagorean Step}

Use one Pythagorean identity, then choose the sign from the range of
$\theta$.

\begin{enumerate}[resume]
\exercise{$3\cos\!\left(\arcsin\dfrac{x}{3}\right)$}{30mm}{%
$\theta\in[-\pi/2,\pi/2]$ so $\cos\theta\geq 0$:
$\cos\theta=\sqrt{1-x^2/9}$ and the value is $\sqrt{9-x^2}$ for
$|x|\leq 3$.}

\exercise{$\sin\!\left(\arccos 4x\right)$}{30mm}{%
$\theta\in[0,\pi]$ so $\sin\theta\geq 0$: the value is
$\sqrt{1-16x^2}$ for $|x|\leq 1/4$.}

\exercise{$\tan\!\left(\arcsin\dfrac{x}{2}\right)$}{32mm}{%
$\sin\theta = x/2$, $\cos\theta = \sqrt{4-x^2}/2>0$, so
$\dfrac{x}{\sqrt{4-x^2}}$ for $|x|<2$.}

\exercise{$5\sec\!\left(\arctan\dfrac{x}{5}\right)$}{32mm}{%
$\sec^2\theta = 1+x^2/25$ and $\sec\theta>0$ on $(-\pi/2,\pi/2)$,
so the value is $\sqrt{25+x^2}$ for all $x$.}

\exercise{$\cos\!\left(\arctan 4x\right)$}{30mm}{%
$\dfrac{1}{\sqrt{1+16x^2}}$ for all $x$.}

\exercise{$\cot\!\left(\arccos\dfrac{x}{6}\right)$}{34mm}{%
$\cos\theta = x/6$, $\sin\theta = \sqrt{36-x^2}/6 \geq 0$, so
$\dfrac{x}{\sqrt{36-x^2}}$ for $|x|<6$. The sign of the answer follows the
sign of $x$ automatically; no absolute value appears.}

\exercise{$\csc\!\left(\arccos 2x\right)$}{30mm}{%
$\dfrac{1}{\sqrt{1-4x^2}}$ for $|x|<1/2$.}

\exercise{Simplify $3\cos\!\left(\arcsin(x/4)\right)$.}{32mm}{%
Let $\theta=\arcsin(x/4)$. Then $\sin\theta=x/4$ and
$\cos\theta\geq0$, so
\[
\boxed{3\cos\!\left(\arcsin(x/4)\right)
=\frac{3\sqrt{16-x^2}}{4}},\qquad |x|\leq4.
\]}

\exercise{Simplify $2\sin\!\left(\arctan(3x)\right)$.}{32mm}{%
Let $\theta=\arctan(3x)$. Since $\tan\theta=3x$ and
$\sec\theta=\sqrt{1+9x^2}$,
\[
\boxed{2\sin\!\left(\arctan(3x)\right)
=\frac{6x}{\sqrt{1+9x^2}}}.
\]}

\exercise{For $0<x\leq\tfrac12$, simplify
$\tfrac15\tan\!\left(\arccos(2x)\right)$.}{34mm}{%
Let $\theta=\arccos(2x)$. Then $\cos\theta=2x>0$ and
$\sin\theta=\sqrt{1-4x^2}$. Thus
\[
\boxed{\frac15\tan\!\left(\arccos(2x)\right)
=\frac{\sqrt{1-4x^2}}{10x}}.
\]}

\exercise{Simplify $\sec\!\left(\arcsin(2x)\right)$.}{32mm}{%
Let $\theta=\arcsin(2x)$. Then
$\cos\theta=\sqrt{1-4x^2}$. Consequently
\[
\boxed{\sec\!\left(\arcsin(2x)\right)
=\frac{1}{\sqrt{1-4x^2}}},\qquad |x|<\frac12.
\]}
\end{enumerate}

\section{Sign-Sensitive Compositions}

These are the cases where the principal range can force an absolute value or
a sign factor.

\begin{enumerate}[resume]
\exercise{$\sin\!\left(\arcsec\dfrac{x}{3}\right)$}{36mm}{%
$\cos\theta = 3/x$ and $\theta\in[0,\pi]$ forces $\sin\theta\geq0$:
$\sin\theta = \sqrt{1-9/x^2} = \dfrac{\sqrt{x^2-9}}{|x|}$ for $|x|\geq 3$.
The naive triangle gives $\sqrt{x^2-9}/x$, which is wrong for $x\leq -3$.}

\exercise{$\tan\!\left(\arcsec\dfrac{x}{2}\right)$}{36mm}{%
$\cos\theta = 2/x$, $\sin\theta = \sqrt{x^2-4}/|x|$, so
\[
\tan\theta = \dfrac{\sin\theta}{\cos\theta}
= \dfrac{x}{|x|}\cdot\dfrac{\sqrt{x^2-4}}{2}
= \sgn(x)\,\dfrac{\sqrt{x^2-4}}{2}
\]
for $|x|>2$.}

\exercise{$\cot\!\left(\arccsc\dfrac{x}{5}\right)$}{36mm}{%
$\sin\theta = 5/x$ and $\theta\in[-\pi/2,\pi/2]$ forces
$\cos\theta\geq 0$: $\cos\theta = \sqrt{x^2-25}/|x|$, so
$\cot\theta = \sgn(x)\,\dfrac{\sqrt{x^2-25}}{5}$ for $|x|>5$.}

\exercise{$\cos\!\left(\arccsc 3x\right)$}{34mm}{%
$\sin\theta = 1/(3x)$, $\cos\theta \geq 0$, so
$\dfrac{\sqrt{9x^2-1}}{3|x|}$ for $|x|\geq 1/3$.}

\exercise{$\sec\!\left(\arccsc\dfrac{x}{2}\right)$}{34mm}{%
From the previous pattern $\cos\theta = \sqrt{x^2-4}/|x|$, so
$\sec\theta = \dfrac{|x|}{\sqrt{x^2-4}}$ for $|x|>2$.}

\exercise{$\csc\!\left(\arcsec 4x\right)$}{36mm}{%
$\cos\theta = 1/(4x)$, $\sin\theta = \sqrt{16x^2-1}/(4|x|)\geq 0$, so
$\dfrac{4|x|}{\sqrt{16x^2-1}}$ for $|x|>1/4$.}

\exercise{$\sin\!\left(\arccot\dfrac{x}{4}\right)$}{34mm}{%
$\cot\theta = x/4$ and $\theta\in(0,\pi)$ forces $\sin\theta>0$:
$\csc^2\theta = 1+x^2/16$ gives $\csc\theta = \sqrt{16+x^2}/4$, so
$\dfrac{4}{\sqrt{16+x^2}}$ for all $x$. The answer is positive even for
$x<0$.}

\exercise{$\cos\!\left(\arccot\dfrac{x}{4}\right)$}{34mm}{%
$\cos\theta = \cot\theta\sin\theta
= \dfrac{x}{4}\cdot\dfrac{4}{\sqrt{16+x^2}}
= \dfrac{x}{\sqrt{16+x^2}}$ for all $x$. Compare with the previous item:
here the sign of $x$ survives, there it does not.}
\end{enumerate}

\section{Double-Angle Compositions}

First recover the inner trigonometric quantities, then apply a double-angle
identity.

\begin{enumerate}[resume]
\exercise{$\sin\!\left(2\arcsin\dfrac{x}{3}\right)$}{38mm}{%
$2\sin\theta\cos\theta
= 2\cdot\dfrac{x}{3}\cdot\dfrac{\sqrt{9-x^2}}{3}
= \dfrac{2x\sqrt{9-x^2}}{9}$ for $|x|\leq 3$.}

\exercise{$\cos\!\left(2\arctan 2x\right)$}{38mm}{%
$\cos 2\theta = \dfrac{1-\tan^2\theta}{1+\tan^2\theta}
= \dfrac{1-4x^2}{1+4x^2}$ for all $x$.}

\exercise{$\sin\!\left(2\arctan\dfrac{x}{5}\right)$}{38mm}{%
$\sin 2\theta = \dfrac{2\tan\theta}{1+\tan^2\theta}
= \dfrac{2x/5}{1+x^2/25} = \dfrac{10x}{25+x^2}$ for all $x$.}

\exercise{$\tan\!\left(2\arctan\dfrac{x}{3}\right)$}{38mm}{%
$\dfrac{2\tan\theta}{1-\tan^2\theta}
= \dfrac{2x/3}{1-x^2/9} = \dfrac{6x}{9-x^2}$ for $x\neq\pm 3$.
The tangent is also undefined where $9-x^2=0$, the same points.}

\exercise{Simplify $2\sin\!\left(2\arctan x\right)$.}{38mm}{%
Put $\theta=\arctan x$. Then
$\sin\theta=x/\sqrt{1+x^2}$ and
$\cos\theta=1/\sqrt{1+x^2}$. Hence
\[
\boxed{2\sin(2\theta)=4\sin\theta\cos\theta
=\frac{4x}{1+x^2}}.
\]}

\exercise{Simplify $\cos\!\left(2\arcsin(x/3)\right)$.}{36mm}{%
Put $\theta=\arcsin(x/3)$. Using
$\cos(2\theta)=1-2\sin^2\theta$ gives
\[
\boxed{\cos\!\left(2\arcsin(x/3)\right)
=1-\frac{2x^2}{9}},\qquad |x|\leq3.
\]}

\exercise{Simplify
$\tfrac12\tan\!\left(2\arctan(x/2)\right)$.}{40mm}{%
Put $\theta=\arctan(x/2)$, so $\tan\theta=x/2$. Then
\[
\frac12\tan(2\theta)
=\frac12\frac{2(x/2)}{1-x^2/4}
=\boxed{\frac{2x}{4-x^2}},\qquad x\neq\pm2.
\]
At $x=\pm2$ the original tangent is also undefined.}
\end{enumerate}

\section{Final Mixed Review}

Use whichever method is appropriate. The section title no longer tells which
identity or sign issue is involved.

\begin{enumerate}[resume]
\exercise{For $x\geq3$, simplify
$4\sin\!\left(\arcsec(x/3)\right)$.}{36mm}{%
Let $\theta=\arcsec(x/3)$. Then $\cos\theta=3/x$ and
$\sin\theta\geq0$. Therefore
\[
\boxed{4\sin\!\left(\arcsec(x/3)\right)
=\frac{4\sqrt{x^2-9}}{x}}.
\]}

\exercise{For $x\neq0$, simplify
$\cot\!\left(\arctan(x/2)\right)$.}{30mm}{%
If $\theta=\arctan(x/2)$, then $\tan\theta=x/2$. Thus
\[
\boxed{\cot\!\left(\arctan(x/2)\right)=\frac{2}{x}}.
\]}

\exercise{$\cos\!\left(2\arcsec\dfrac{x}{2}\right)$}{38mm}{%
$2\cos^2\theta - 1 = \dfrac{8}{x^2}-1 = \dfrac{8-x^2}{x^2}$ for
$|x|\geq 2$.}

\exercise{$\sin\!\left(2\arcsec\dfrac{x}{4}\right)$}{40mm}{%
$2\sin\theta\cos\theta
= 2\cdot\dfrac{\sqrt{x^2-16}}{|x|}\cdot\dfrac{4}{x}
= \dfrac{8\sgn(x)\sqrt{x^2-16}}{x^2}$ for $|x|\geq 4$.}

\exercise{For $x\geq2$, simplify
$3\cos\!\left(2\arcsec(x/2)\right)$.}{38mm}{%
Put $\theta=\arcsec(x/2)$, so $\cos\theta=2/x$.
Using $\cos(2\theta)=2\cos^2\theta-1$,
\[
\boxed{3\cos\!\left(2\arcsec(x/2)\right)
=3\left(\frac{8}{x^2}-1\right)}.
\]}

\exercise{Simplify $\sin\!\left(2\arccos x\right)$.}{40mm}{%
Put $\theta=\arccos x$. Since $\theta\in[0,\pi]$,
$\sin\theta=\sqrt{1-x^2}$, while $\cos\theta=x$. Thus
\[
\boxed{\sin\!\left(2\arccos x\right)
=2x\sqrt{1-x^2}},\qquad |x|\leq1.
\]
The factor $x$ correctly supplies the sign when $\theta>\pi/2$.}

\end{enumerate}

\section*{Final check}

For each answer, ask:
\begin{enumerate}[label=\arabic*.]
  \item Is the answer entirely in terms of $x$?
  \item Did I use the correct Pythagorean, reciprocal, or double-angle identity?
  \item Did I choose the sign from the inverse function's principal range or
        the stated range of $\theta$?
  \item Does the final expression have the same domain as the original one?
\end{enumerate}

\end{document}
